% Part: first-order-logic
% Chapter: completeness
% Section: downward-ls

\documentclass[../../../include/open-logic-section]{subfiles}

\begin{document}

\olfileid[pt]{fol}{com}{dls}
\olsection{O Teorema de L\"owenheim--Skolem}

O Teorema de L\"owenheim--Skolem diz que se uma teoria tem um modelo
infinito, então ela também tem um modelo que é no máximo
!!{denumerable}. Uma consequência imediata desse fato é que a lógica de
primeira ordem não pode expressar que o tamanho de !!a{structure} é
!!{nonenumerable}: qualquer !!{sentence} ou conjunto de !!{sentence}s
satisfeito em todas as !!{structure}s !!{nonenumerable}s também é
satisfeito em alguma estrutura !!{enumerable}.

\begin{thm}
\ollabel{thm:downward-ls} Se $\Gamma$ é consistente, então ela tem
um modelo !!{enumerable}, isto é, é satisfatível em !!a{structure}
cujo domínio é finito ou !!{denumerable}.
\end{thm}

\begin{proof}
Se $\Gamma$ é consistente, a !!{structure}~$\Struct M$ fornecida pela
prova do teorema da completude tem um domínio $\Domain{M}$ que não é
maior do que o conjunto dos termos da linguagem~$\Lang L$. Assim
$\Struct M$ é no máximo !!{denumerable}.
\end{proof}

\begin{thm}
\ollabel{noidentity-ls} Se $\Gamma$ é um conjunto consistente de
!!{sentence}s na linguagem da lógica de primeira ordem sem identidade,
então ela tem um modelo !!{denumerable}, isto é, é satisfatível em
!!a{structure} cujo domínio é infinito e !!{enumerable}.
\end{thm}

\begin{proof}
Se $\Gamma$ é consistente e não contém sentenças em que a identidade
aparece, então a !!{structure}~$\Struct M$ fornecida pela prova do
teorema da completude tem um domínio $\Domain{M}$ idêntico ao conjunto
dos termos da linguagem~$\Lang L'$. Assim $\Struct{M}$ é
!!{denumerable}, pois $\Trm[L']$ o é.
\end{proof}

\begin{ex}[Paradoxo de Skolem]
A teoria dos conjuntos de Zermelo--Fraenkel~$\Log{ZFC}$ é um arcabouço
muito poderoso no qual praticamente todas as afirmações matemáticas
podem ser expressas, incluindo fatos sobre os tamanhos de conjuntos.
Assim, por exemplo, $\Log{ZFC}$ pode provar que o conjunto~$\Real$ dos
números reais é !!{nonenumerable}, pode provar o Teorema de Cantor de
que o conjunto potência de qualquer conjunto é maior do que o próprio
conjunto, etc.  Se $\Log{ZFC}$ é consistente, seus modelos são todos
infinitos e, além disso, todos contêm !!{element}s sobre os quais a
teoria diz que são !!{nonenumerable}s, como o elemento que torna
verdadeiro o teorema de~$\Log{ZFC}$ de que o conjunto potência dos
números naturais existe. Pelo Teorema de L\"owenheim--Skolem,
$\Log{ZFC}$ também tem modelos !!{enumerable}s---modelos que contêm
conjuntos ``!!{nonenumerable}s'', mas que eles mesmos são
!!{enumerable}s.
\end{ex}

\end{document}
